% SGA 3，第 VII 篇_A，中文重构——第 1 轮。
% 范围：权威底本本地第 27--28 页，4.2--4.2.1，止于本地第 29 页的 4.3 之前。
% 合并阅读本页码： 480--481. 再版本印刷页码： 469--470.
% 本段原版页边定位： 440--441.
% 权威底本：Polo--Gille Exp7A-13oct24.pdf,
% SHA-256 D46DB8C79D0385643B6EED4B96900613248375F9B73B4202478A4F57FD3E8522.
% 参照了 Jacob Reinhold 的 CC BY 4.0 Markdown 作为比对谱系。

\medskip
\noindent\textbf{4.2.}
\phantomsection
\label{sga3:VIIA:4.2}
\footnote{关于 \SGARefLink{sga3:VIIA:4.2}{4.2} 及 \SGARefLink{sga3:VIIA:4.3}{4.3} 的内容，另见
\textup{[DG70]}, IV, \S3, 第~4--6~目.}

% 权威底本： Exp7A 本地页 27; 合并阅读本页 480; 印刷页 469.
现在考虑一个与前述构造密切相关的构造。设 \(S\) 是特征为 \(p\) 的概形，\(X\) 是一个 \(S\)-概形，\(X_S^p\) 是范畴 \((\mathbf{Sch}/S)\) 中 \(p\) 个 \(X\) 的乘积。

以 \(U^p(X)\) 表示 \(X_S^p\) 的开子概形，它是当 \(U\) 遍历 \(X\) 的仿射开子集时各乘积 \(U_S^p\) 的并。因此，\(X_S^p\) 的点 \(x\) 属于 \(U^p(X)\)，当且仅当它在 \(X_S^p\) 的各因子上的投影 \(\operatorname{pr}_i x\) 属于 \(X\) 的同一个仿射开子集。例如，若 \(X\) 的每个有限子集都包含在一个仿射开子集中，则
\[
  U^p(X)=X_S^p.
\]

\(p\) 个字母上的对称群 \(\mathfrak S_p\) 通过置换因子作用于 \(X_S^p\)，并保持 \(U^p(X)\) 稳定。我们把带环空间范畴中 \(X_S^p\) 关于 \(\mathfrak S_p\) 的商称为 \emph{\(X\) 的 \(p\) 重对称积}，记作 \(\Sigma^p X\)。设
\[
  q(X),\ \text{或简记为 }q:
  X_S^p\longrightarrow\Sigma^p X
\]
为典范投影。

态射 \(q\) 把 \(U^p(X)\) 映满到对称积的一个开子集 \(V^p(X)\)，其描述如下（参见 V，\SGARefLink{thm:sga3-V-4.1}{4.1}）。\(\Sigma^p X\) 的结构层在 \(V^p(X)\) 上诱导一个概形结构，而由 \(q(X)\) 诱导的态射
\[
  q'(X):U^p(X)\longrightarrow V^p(X)
\]
是仿射的，而且是整的。当 \(U\) 遍历映入 \(S\) 的某个变动的仿射开集 \(V\) 的 \(X\) 的仿射开子集时，概形 \(\Sigma^p U\) 构成 \(V^p(X)\) 的一个仿射开覆盖。若 \(R\) 是 \(V\) 的仿射代数，\(A\) 是 \(U\) 的仿射代数，则 \(\Sigma^p U\) 的仿射代数是 \(\bigotimes_R^p A\) 中由对称张量组成的子代数 \(\Sigma^p A\)。

现在考虑对角态射
\[
  \delta:X\longrightarrow U^p(X).
\]
若 \(V=\operatorname{Spec}R\) 是 \(S\) 的仿射开子集，\(U=\operatorname{Spec}A\) 是位于 \(V\) 之上的 \(X\) 的仿射开子集，则 \(\delta\) 在 \(U\) 上的限制由如下代数同态定义
\[
  \eta:\bigotimes_R^p A\longrightarrow A,
  \qquad
  a_1\otimes\cdots\otimes a_p
  \longmapsto
  a_1a_2\cdots a_p.
\]
因此，若 \(N\) 表示对称化算子，则
\[
\begin{aligned}
  \eta\bigl(N(a_1\otimes\cdots\otimes a_p)\bigr)
  &=
  \eta\left(
    \sum_{\sigma\in\mathfrak S_p}
      a_{\sigma(1)}\otimes\cdots\otimes a_{\sigma(p)}
  \right) \\
  &=
  p!\,a_1\cdots a_p
  =
  0.
\end{aligned}
\]
换言之，\(\eta\) 在由对称化张量组成的子空间
\[
  N\!\left(\bigotimes_R^p A\right)
  \subseteq
  \Sigma^p A
\]
上为零。此外，若 \(f\) 是对称张量，则显然 \(N(fa)=fN(a)\)。故 \(N(\bigotimes_R^p A)\) 是 \(\Sigma^p A\) 的一个理想。置
\[
  U^{[p/S]}
  =
  \operatorname{Spec}\left(
    \Sigma^p A
    \big/
    N\!\left(\bigotimes_R^p A\right)
  \right).
\]
这是 \(\Sigma^p(U)=V^p(U)\) 的一个闭子概形。当 \(U\) 遍历映入 \(S\) 的某个变动的仿射开集 \(V\) 的 \(X\) 的仿射开子集时，概形 \(U^{[p/S]}\) 粘合成 \(V^p(X)\) 的一个闭子概形，记作 \(X^{[p/S]}\)。

若 \(i(X)\) 表示 \(X^{[p/S]}\) 到 \(V^p(X)\) 的包含态射，则前面的计算表明 \(q'(X)\circ\delta\) 经由 \(X^{[p/S]}\) 分解。因此得到态射
\[
  F^{[p]}(X/S):X\longrightarrow X^{[p/S]}.
\]
\footnote{原版中，此态射和相对 Frobenius 态射分别记作 \(F'\) 和 \(F\)。}
此分解由如下交换图表示
\begin{center}
  \phantomsection
  \label{sga3:VIIA:4.2-pfold-symmetric-factorization}
  % SGA3-VII-DIAG-044
  \begin{tikzcd}[row sep=3.2em,column sep=1.8em]
X_S^p \arrow[dd,swap,"q(X)"]
  & \supset
  & \mathrm U^p(X) \arrow[dd,swap,"q'(X)"]
  && X \arrow[ll,swap,"\delta(X)"] \arrow[dd,"F^{[p]}(X/S)"] \\
\\[-1.2em]
\Sigma^p(X)
  & \supset
  & \mathrm V^p(X)
  && X^{[p/S]} . \arrow[ll,swap,"i(X)"]
\end{tikzcd}

\end{center}
显然，\(X^{[p/S]}\) 关于 \(X\) 具有函子性，且
\[
  F^{[p]}:X\longmapsto F^{[p]}(X/S)
\]
是一个函子性同态。

\medskip
\noindent\textbf{4.2.1.}
\phantomsection
\label{sga3:VIIA:4.2.1}

% 权威底本： Exp7A 本地页 28; 合并阅读本页 481; 印刷页 470.
概形 \(X^{[p/S]}\) 与 \(X^{(p/S)}\) 有自然的联系。设 \(V\) 是 \(S\) 的仿射开子集，其仿射代数为 \(R\)；\(U\) 是位于 \(V\) 之上的 \(X\) 的仿射开子集，其仿射代数为 \(A\)。若 \(\pi\) 表示 \(R\) 的自同态 \(x\mapsto x^p\)，则 \(U^{(p/S)}\) 的仿射代数为 \(A\otimes_\pi R\)。映射
\[
  a\otimes_\pi\lambda
  \longmapsto
  \left(
    \lambda a\otimes\cdots\otimes a
    \bmod
    N\!\left(\bigotimes_R^p A\right)
  \right)
\]
是一个 \(R\)-代数同态
\[
  A\otimes_\pi R
  \longrightarrow
  \Sigma^p A\big/
    N\!\left(\bigotimes_R^p A\right).
\]
它诱导态射
\[
  \varphi(U):U^{[p/S]}\longrightarrow U^{(p/S)}
\]
满足
\[
  \varphi(U)\circ F^{[p]}(U/S)=\operatorname{Fr}(U/S).
\]
粘合这些态射，得到如下交换三角图
\begin{center}
  \phantomsection
  \label{sga3:VIIA:4.2.1-phi-frobenius-triangle}
  % SGA3-VII-DIAG-045
  \begin{tikzcd}[row sep=3.6em,column sep=4.4em]
& X
  \arrow[dl,swap,"F^{[p]}(X/S)"]
  \arrow[dr,"\operatorname{Fr}(X/S)"] & \\
X^{[p/S]} \arrow[rr,"\varphi(X)"]
  && X^{(p/S)} .
\end{tikzcd}

\end{center}

例如，若 \(X\) 是由拟凝聚理想 \(\mathcal I\) 定义的 \(S\) 的子概形，则 \(F^{[p]}(X/S)\) 可等同于 \(X\) 的恒等态射。因此 \(\varphi(X)\) 是典范浸入
\[
  \operatorname{Spec}(\mathcal O_S/\mathcal I)
  \longrightarrow
  \operatorname{Spec}(\mathcal O_S/\mathcal I^{\{p\}}).
\]
这表明 \(\varphi(X)\) 未必是同构。

另一方面，当 \(M\) 是自由 \(R\)-模时，映射
\[
\begin{aligned}
  M\otimes_\pi R
  &\longrightarrow
  \Sigma^p M\big/
    N\!\left(\bigotimes_R^p M\right), \\
  m\otimes_\pi\lambda
  &\longmapsto
  \left(
    \lambda m\otimes\cdots\otimes m
    \bmod
    N\!\left(\bigotimes_R^p M\right)
  \right)
\end{aligned}
\]
显然是双射。当 \(M\) 在 \(R\) 上平坦时，它仍是双射，因为每个平坦模都是自由模的滤过归纳极限。\footnote{D.~Lazard,
\emph{关于平坦模〔Sur les modules plats〕}, C.\ R.\ Acad.\ Sci.\ Paris
\textbf{258} (1964), 6313--6316. 另见
D.~Lazard, Bull.\ Soc.\ Math.\ France \textbf{97} (1969), 81--128,
或 \textup{[BAlg]}, X, \S1.6, 定理~1.}
于是
\[
  \boxed{\;
  \varphi(X):X^{[p/S]}\longrightarrow X^{(p/S)}
  \text{ 在 \(X\) 于 \(S\) 上平坦时是同构。}
  \;}
\]

% 下一定位： Exp7A 本地页 29, 合并阅读本页 482,
% 再版本印刷页 471, 在 4.3。
